\begin{lemma}[The technical theorem at rung zero]
For the initial tuple
((T_1,\varphi_0,S_0,F_0,\Theta_0)) of every certified Tashkinov series,
the rung-zero package \noderef{TechnicalTheoremBaseClaims} holds: (T_1)
is elementary under (arphi_0), and (T_1) has the interchangeability
property under every proper ((T_1,D_0,\varphi_0))-stable coloring leaving
exactly (e_0) uncolored.  The indexed positive-rung assertions are
vacuous.
\end{lemma}
\begin{proof}
Write (T_0,T_1) and (arphi_0) for data slots (0,1) of the
series.  Eliminating the certificate at slot (0) in
\noderef{TashkinovSeries} leaves precisely the initial constructor of
\noderef{TashkinovSeriesStep}: it states that
((G,e_0,\varphi_0)) is a (k)-triple, that (T_0) has length one and
first edge (e_0), that (S_0=F_0=\varnothing), and that (T_1) is a
closure of (T_0) under (arphi_0).  By
\noderef{TashkinovStoredColors}, (D_0=\varnothing).

The length-one tree (T_0) is a Tashkinov tree under (arphi_0): in
the defining condition \noderef{TashkinovTree}, the only edge position is
position zero, while the condition asks only about positions at least one.
The closure assertion, by \noderef{TashkinovClosure}, supplies a
\noderef{TashkinovAugmentationSequence} from (T_0) to (T_1) and says
that (V(T_1)) is closed under (arphi_0).  Applying
\noderef{TashkinovAugmentationSequencePreservesTashkinovTree} gives that
(T_1) is a Tashkinov tree under (arphi_0).  The (k)-triple and this
last fact are exactly the hypotheses of \noderef{TashTree}; hence
(V(T_1)) is elementary under (arphi_0).  This proves the first field
of the base package.

Now let (sigma_0) be proper, leave precisely (e_0) uncolored, and be
((T_1,D_0,\varphi_0))-stable.  We first verify that the same displayed
(T_1) is a closed Tashkinov tree under (sigma_0).

Consider one augmentation (A\to B) in the preceding augmentation
sequence.  By \noderef{TashkinovAugmentation}, its new edge (f) has
(arphi_0(f)=c), where (c) is missing at an old vertex (x\in V(A)).
Every augmentation preserves the previous display as the ordered initial
segment specified by \noderef{TreeSequenceInitialSegment}; induction along
the finite sequence therefore puts (x) in (V(T_1)) and makes (f)
incident with (T_1).  Stability in the sense of
\noderef{StableColoring} gives
(overline\sigma_0(x)=\overline\varphi_0(x)), and its incident-edge
clause fixes (f), because (c\in\overline\varphi_0(x)\subseteq
\overline\varphi_0(T_1)).  Thus
(sigma_0(f)=\varphi_0(f)=c), with (c) missing at (x) also under
(sigma_0).  Consequently the same step is a Tashkinov augmentation
under (sigma_0).  Induction over the relation
\noderef{TashkinovAugmentationSequence} transports the entire sequence
from (T_0) to (T_1).  The length-one argument makes (T_0) a
Tashkinov tree under (sigma_0), so
\noderef{TashkinovAugmentationSequencePreservesTashkinovTree} now makes
(T_1) a Tashkinov tree under (sigma_0).

For closedness, suppose that a boundary edge (f) had
(sigma_0(f)=c), where (c) is missing at a vertex of (T_1) under
(sigma_0).  Stability identifies that missing set with the one under
(arphi_0).  Both colorings are proper and leave the same unique edge
uncolored, so the reverse-color clause of \noderef{sc1} applies to the
incident boundary edge and gives
(arphi_0(f)=sigma_0(f)=c).  This contradicts the closedness of
(V(T_1)) under (arphi_0) supplied by the initial closure and defined
in \noderef{ColoredVertexSetClosed}.  Hence (T_1) is closed under
(sigma_0).

The graph-theoretic, chromatic-index, and degree clauses of
\noderef{KTriple} are independent of the coloring.  Taking those clauses
from the initial (k)-triple and adjoining the assumed properness and
unique-uncolored properties of (sigma_0) shows that
((G,e_0,\sigma_0)) is a (k)-triple.  We have proved all three remaining
hypotheses of \noderef{Tchange}: this (k)-triple, the Tashkinov-tree
property of (T_1), and its closedness under (sigma_0).  Therefore, for
any colors (alpha,eta) of which at least one is missing on (T_1),
they are (T_1)-interchangeable under (sigma_0).  This is exactly the
second field of \noderef{TechnicalTheoremBaseClaims}, and the proof is
complete.
\end{proof}
